\subsection{Fusing the first vote with the bias correction}
\label{app:fused}

For row norms at most one, the uniform algorithm uses reference gain
one in the radius recurrence. The exact row norms still determine the row
mixtures and their radii \(\rho\). This choice preserves the sampler,
its exactness proof, and the bound \(80(D+1)^2\). We prove a smaller
power of depth in two steps. First we assign a positive weight to each
reached call; the weights are used only in the analysis. Then, at
sufficiently small row radius, we replace the bias race by a factory
whose first vote also selects the bias correction. This reduces the
critical power without changing the represented function.

\subsubsection{Weighted local estimates for the critical sampler}
\label{app:potential}

Fix \(0<\eta\le1/4\), write \(\alpha=1+\eta\), and define
\[
\varphi(v)=\sqrt{\alpha-v^2},\qquad
\ell(v)=\log\varphi(v),\qquad -1\le v\le1.
\]
The function \(\varphi\) is even and concave, with
\(\sqrt\eta\le\varphi\le\sqrt\alpha\). Direct differentiation gives
\begin{equation}\label{eq:potentialderivatives}
\ell'(v)=-\frac{v}{\alpha-v^2},\qquad
\ell''(v)=-\frac{\alpha+v^2}{(\alpha-v^2)^2},\qquad
|\ell'|\le\eta^{-1},\quad
|\ell''|\le\frac9{4\eta^2}.
\end{equation}

At a nonzero row put
\[
q=\tanh\rho,\quad g=\frac{\tanh(\rho z)}q,\quad t=\tanh a.
\]
If \(t<0\), complement source and output signs; evenness of
\(\varphi\) leaves the analysis unchanged. Assume henceforth that
\(t\ge0\), and set
\[
c=tq,\qquad k=\frac{1-t^2}{1-c^2},\qquad
T_c(g)=\frac{g+c}{1+cg}.
\]
The exact midpoint and radius satisfy
\(R=qk\) and \(h=m+RT_c(g)\). For the stored center
\(c_{\rm row}\) and radius \(r\), let
\[
\theta=q/r\le1,\qquad \delta=(m-c_{\rm row})/r.
\]
The parent's normalized mean is
\(v=\delta+\theta kT_c(g)\in[-1,1]\).

Assign a child of normalized mean \(v_j\) the analysis weight
\(\varphi(v_j)\). If its row-mixture probability and sign are
\(p_j\) and \(\sigma_j\), Jensen's inequality gives
\[
\sum_jp_j\varphi(v_j)
=\sum_jp_j\varphi(\sigma_jv_j)
\le\varphi\!\left(\sum_jp_j\sigma_jv_j\right)=\varphi(z).
\]
Applying Lemma~\ref{lem:stopping} first to the majority and then
to the race therefore gives
\begin{equation}\label{eq:weightedchildren}
\E\sum_{\rm children}\varphi(v_j)
\le\theta kA_{\rm seq}(\rho,z)\frac{1+c}{1+cg}\varphi(z).
\end{equation}
This use of the lemma permits correlation between a source's cost
and returned sign. The event that requests a source is determined
before its fresh randomness is used.

For fixed \(|u|\le1\), the function
\(\theta/\varphi(\theta u)\) is increasing on \([0,1]\).
Both \(\theta kT_c(g)\) and \(v\) lie in \([-1,1]\), so
\eqref{eq:potentialderivatives} bounds the effect of the center
shift by \(\exp(|\delta|/\eta)\). Consequently the right side
of \eqref{eq:weightedchildren}, divided by \(\varphi(v)\), is at most
\begin{equation}\label{eq:potentialW}
\exp\!\left(\frac\varepsilon{\eta r}\right)W,
\qquad
W=A_{\rm seq}(\rho,z)k\frac{1+c}{1+cg}
\frac{\varphi(z)}{\varphi(kT_c(g))}.
\end{equation}

\begin{lemma}[Majority cost and the bias contribution]
\label{lem:potential-majority-bias}
For $0<\eta\le1/4$ and $0<\rho\le1/16$, the majority term
satisfies
\[
 \log A_{\rm seq}+\ell(z)-\ell(g)
 \le\left(1-\frac{\eta z^2}{3(\alpha-z^2)}\right)\rho^2
      +\frac4{\eta^2}\rho^4.
\]
The remaining bias term
$B=\log k+\log((1+c)/(1+cg))+\ell(g)-\ell(kT_c(g))$
obeys
\[
 B\le\frac{-(1-2q^2)t^2+3tq}{1-q^2}.
\]
In particular $B\le0$ when $t\ge(7/2)q$.
\end{lemma}

\paragraph{Proof idea.}
A real Taylor expansion controls the majority contribution. For the bias
term, express the potential along the coordinate $\operatorname{arctanh}g$.
Its derivative is bounded by one, which controls the whole translation even
near the endpoints.

\begin{proof}
First record Taylor bounds that hold throughout \(\rho\le1/16\):
\begin{equation}\label{eq:potentialgtaylor}
\left|g-z-\frac{\rho^2z(1-z^2)}3\right|\le11\rho^4,
\qquad |g-z|\le\rho^2/2.
\end{equation}
Indeed, for complex \(|w|\le1\),
\(|\cosh w-1|\le\cosh1-1<3/5\) and
\(|\sinh w|\le\sinh1<6/5\), by \(8/3<e<11/4\); hence
\(|\tanh w|<3\). Cauchy's coefficient estimate and oddness give
\begin{equation}\label{eq:tanhfifth}
|\tanh w-w+w^3/3|\le5|w|^5\qquad(w\in\mathbb C,\ |w|\le1/16).
\end{equation}
Also
\(q/\rho\ge1-\rho^2/3\ge767/768\), by integrating
\(\tanh'(u)\ge1-u^2\). Subtracting \(zq/\rho\) from the
numerator expansion leaves an error at most \(10\rho^4\);
dividing by \(q/\rho\) adds less than \(\rho^4\).
Finally, \(1/3+11/256<1/2\).

Lemma~\ref{lem:arity}, Taylor's theorem for \(\ell\), and
\eqref{eq:potentialgtaylor} give, with \(d_z=\alpha-z^2\),
\begin{align}
\log A_{\rm seq}+\ell(z)-\ell(g)
&\le\left(1-\frac{\eta z^2}{3d_z}\right)\rho^2
 +\left(1+\frac{11}\eta+\frac9{32\eta^2}\right)\rho^4\notag\\
&\le\left(1-\frac{\eta z^2}{3d_z}\right)\rho^2
 +\frac4{\eta^2}\rho^4.\label{eq:potentialmajority}
\end{align}
The second line uses
\(\eta^2+11\eta+9/32\le99/32<4\).

It remains to bound the bias contribution
\[
B=\log k+\log\frac{1+c}{1+cg}+\ell(g)-\ell(kT_c(g)).
\]
Since \(k\le1\), the last difference
is at most \(\ell(g)-\ell(T_c(g))\). For \(v=\tanh u\),
\[
\left|\frac{d}{du}\ell(\tanh u)\right|
=\frac{|v|(1-v^2)}{\alpha-v^2}\le1.
\]
Thus that difference is at most \(\artanh c\), including
\(g=\pm1\) by continuity. The race logarithm is at most
\(2\artanh c\). Since
\[
\log k\le-\frac{1-2q^2}{1-q^2}t^2,
\qquad \artanh c\le\frac{tq}{1-q^2},
\]
we obtain
\[
B\le\frac{-(1-2q^2)t^2+3tq}{1-q^2}.
\]
If \(t\ge(7/2)q\), its numerator is at most
\(tq(-1/2+7q^2)\le0\), because \(q\le\rho\le1/16\).

\end{proof}

\paragraph{Bounds at a fixed cutoff.}
For gain one, $\coth^2u=1+\sinh^{-2}u\le1+u^{-2}$ gives
$R_j^2\ge1/(j+1)$.  Lemmas~\ref{lem:radius} and~\ref{lem:rounding}
therefore give
\[
 (j+1)^{-1}\le R_j^2\le3/(2j+3),\quad
 0\le r_j-R_j\le3j\varepsilon,\quad r_j\le1.
\]
The original normalized sampler at level $J$ obeys
\begin{equation}\label{eq:normalizedcutoff}
 M_J\le1+e^3(J+1)^{5/2}\sum_{k\ge1}k^{-5/2}
 \le36(J+1)^{5/2}.
\end{equation}
This follows from the suffix proof of Theorem~\ref{thm:upper},
with the root gate omitted, and $\sum k^{-5/2}\le5/3$, $e^3<21$.
For $\eta=1/64$, $J\ge2^{25}$, $D>J$, the center terms in
\eqref{eq:potentialW} sum to at most
$64\varepsilon D\sqrt{D+1}<1/100$. If a squared-radius potential
has derivative at most two, the proof of Lemma~\ref{lem:rounding}
shows that rounding changes a suffix logarithm
by at most $12D^2\varepsilon<1/100$. These estimates use
$\varepsilon\le[100(D+1)^4]^{-1}$.
